\chapter{Pistes d'amélioration, chapitre par chapitre}\label{annexe:pistes}
\iterbanner{V1}

Le chapitre~\ref{ch:validation} donne les pistes globales. Cette annexe les détaille partie par partie, de la plus utile à la moins utile pour le hackathon.

\section*{Géométrie (ch.~\ref{ch:naca})}
\begin{itemize}
\item Ajouter une paramétrisation qui découple le rayon de bord d'attaque de l'épaisseur (NACA 4 chiffres modifiée, CST de Kulfan) : c'est la condition pour que traînée et SER entrent vraiment en conflit (cours d'optimisation, chapitre sur le front exact).
\item Prévoir un bord de fuite d'épaisseur finie en option, plus réaliste pour la fabrication.
\end{itemize}

\section*{Panneaux (ch.~\ref{ch:potentiel} et \ref{ch:panneaux})}
\begin{itemize}
\item Passer à des tourbillons linéaires par panneau (formulation d'XFOIL) : erreur en $1/N^2$ au lieu de $1/N$.
\item Ajouter une correction de compressibilité de Kármán-Tsien sur $C_p$, pour un drone à Mach 0,5 à 0,6.
\end{itemize}

\section*{Couche limite (ch.~\ref{ch:cl} à \ref{ch:head})}
\begin{itemize}
\item Remplacer Michel par un critère qui tient compte du gradient de pression, ou par une méthode $e^N$ simplifiée (enveloppe de Drela et Giles), et caler sur XFOIL.
\item Remplacer Head par la méthode à décalage de Green (\emph{lag-entrainment}), plus juste en gradient adverse.
\item Option « transition forcée à 5~\% » : cas tout turbulent, plus régulier, donc plus favorable au modèle quadratique de la phase~4.
\end{itemize}

\section*{Traînée (ch.~\ref{ch:trainee})}
\begin{itemize}
\item Remplacer la coupure \code{x\_stop} par une itération de couplage avec vitesse de transpiration, qui efface physiquement le point d'arrêt du bord de fuite.
\end{itemize}

\section*{Assemblage, espace de conception (ch.~\ref{ch:modele} et \ref{ch:espace})}
\begin{itemize}
\item \iterbadge{V1}~\textbf{Fait} : intégration avec un modèle SER simplifié et une chaîne d'optimisation complète (\code{Code/optim/}).
\item Paralléliser l'énumération (512 profils indépendants) pour un passage à des grilles plus fines.
\end{itemize}

\section*{Validation (ch.~\ref{ch:validation})}
\begin{itemize}
\item \iterbadge{V1}~\textbf{Fait} : XFOIL installé localement (\code{Code/.tools/}), test V5 étendu à toute la grille (§\ref{sec:v5}). Reste à caler la transition sur XFOIL.
\item Corriger aussi la portance (décambrage visqueux) : notre $C_L$ est surestimé de 25~\% en moyenne à $2^\circ$, ce qui rend la contrainte $C_L\ge0{,}3$ trop permissive (392 profils faisables contre 359 avec XFOIL).
\end{itemize}
